\documentclass[11pt]{article}
\usepackage[letterpaper,hmargin=1in,vmargin=0.88in]{geometry}
\usepackage[T1]{fontenc}
\usepackage{lmodern,microtype,amsmath,amssymb,amsthm,needspace}
\usepackage[hidelinks,pdftitle={Verifying a matrix transpose on tapes},pdfauthor={Charlie Cruz}]{hyperref}
\newtheorem{theorem}{Theorem}[section]
\newtheorem{lemma}[theorem]{Lemma}
\newtheorem{corollary}[theorem]{Corollary}
\newcommand{\TrLang}{L_{\mathrm{tr}}}
\newcommand{\NT}{\mathsf{NTIME}}
\newcommand{\coNT}{\mathsf{coNTIME}}
\newcommand{\coRT}{\mathsf{coRTIME}}
\newcommand{\ENP}{\mathsf E^{\mathsf{NP}}}
\title{Verifying a Matrix Transpose on Tapes}
\author{Charlie Cruz}
\date{25 November 2026}
\begin{document}
\maketitle

\begin{abstract}
Two row-major bit matrices can be checked for transposition in randomized
linear time on tapes. We show how a sound nondeterministic verifier yields
a rational tensor network for partial transposition. Its coefficients count
accepting paths, so the transfer permits arbitrary ambiguity. Under the
partial-transposition theorem of the integer-multiplication manuscript and
Weak ETH, every such verifier takes $\Omega(n\sqrt{\log n})$ time at every
large valid dimension. This would separate nondeterministic and
co-nondeterministic linear time. A second deduction, using the geometric
case split and an inverse sparse-SAT rate, gives either the same verifier
lower bound or an almost-everywhere $\Omega(n^{1+\varepsilon})$ circuit lower
bound in $\ENP$. We reconstruct the partial-transposition estimate from
two inequalities in the pinned release. Keeping the dependence on the
number of correct inputs gives a further conclusion under the same
assumptions. Below this time bound, a sound machine has a short accepting
path on at most a $2^{-(1-1/\sqrt2-o(1))N^2}$ fraction of true pairs
of side length $N$. A row-entropy version also bounds this fraction
within prescribed products of row sets.
\end{abstract}

\section{The decision problem}
Let $k\ge1$, $N=2^k$, and $m_N=k+1+2N^2$. On binary input $1^k0ab$,
interpret $a,b$ as row-major lists of the entries of $N\times N$ bit
matrices $A,B$. Define $\TrLang$ to consist of these strings with
$|a|=|b|=N^2$ and $B=A^\top$. Reject every malformed string. Machines
have a fixed number of tapes, finite alphabets and heads that move at
most one cell per step. Time bounds hold on all branches. Nonuniform
Boolean circuits have unrestricted wiring.

\begin{lemma}\label{lem:upper}
$\TrLang\in\mathsf L\cap\coNT[O(n)]\cap\coRT[O(n)]$, and its
$n$-bit slice has a fan-in-two circuit of size $O(n)$.
\end{lemma}
\begin{proof}
Count the input length $n$ and unary prefix length $k$ in binary.
Amortized carries take linear total tape time. Reject impossible $k$
by comparison with $\log n$ before forming $2^{2k}$, then check the
exact length and delimiter. These checks also use $O(\log n)$ space.
Logspace suffices to compare all pairs $A_{ij},B_{ji}$ by rescanning.
For the complement, accept malformed strings and otherwise guess the
two $k$-bit coordinates of a mismatch. The two offsets are obtained
by concatenating indices in opposite orders and adding the starts of
the matrices. Sequential scans with decrementing counters take $O(n)$
time in total.

For the randomized test choose independent uniform $r,c\in\mathbb F_2^N$.
While scanning $A$ compute $r^\top Ac$ by scanning and rewinding $c$ once
per row and advancing in $r$ once. Rewind both vectors and use them in
opposite roles on $B$. This takes $O(N^2)$ steps, including rewinds.
Transpose pairs always pass. If $D=A-B^\top\ne0$, then $Dc\ne0$ with
probability at least $1/2$, and conditionally $r^\top Dc=1$ with
probability $1/2$. Reject on any detected mismatch in four independent
trials. The false acceptance probability is at most $(3/4)^4<1/3$.

At length $m_N$ the prefix is fixed. A circuit checks it and conjoins
$N^2$ equality gates on the prescribed pairs. At every other length
use the constant zero circuit.
\end{proof}

\section{From accepting paths to a partial transpose}
\paragraph{Geometric premise.}
Fix $D_0$ and $0<\rho<1-1/\sqrt2$. Consider a rational tensor network
of degree at most $D_0$ in bundles of width at most $D_0N$ bits, with
distinct terminals for all $N$ input rows and $N$ output columns. If
on a fraction at least $2^{-\rho N^2}$ of the row assignments $A$ its
output is a nonzero multiple of the basis vector $e_{A^\top}$, the
network has $\Omega_{D_0,\rho}(N\sqrt{\log N})$ vertices. The source
states this estimate assuming Weak ETH~\cite{partial}.
Section~\ref{sec:audit} derives it from graph placement under the
characteristic-zero no-ports hypothesis and an entropy inequality for
paired terminals. Weak ETH enters through the implication to no ports.
Here Weak ETH excludes deterministic $2^{O(v^{1-\delta})}
\operatorname{poly}(L)$ algorithms for 3-SAT for every fixed
$0<\delta<1$, in $v$ variables and encoding length $L$.

\begin{lemma}[Verifier network]\label{lem:network}
Suppose a fixed nondeterministic machine decides $\TrLang$. At a valid
dimension let $T$ bound its time on all inputs of length $m_N$.
There is a rational network with $O(N+\lceil T/N\rceil)$ vertices,
constant bundle degree, and bundle width $O(N)$, whose coefficients
$K(A,B)$ vanish for $B\ne A^\top$. For at least a
$2^{-O(1+\lceil T/N\rceil)}$ fraction of $A$, $K(A,A^\top)>0$.
\end{lemma}
\begin{proof}
Pad halting paths to $T$ with one stationary continuation per step.
Partition time into epochs of at most $N$ steps and each tape into
blocks of $N$ cells. At each epoch boundary record each head's block.
Its initial block is fixed and it moves by at most one block per epoch.
Thus there are at most $3^{h\lceil T/N\rceil}$ schedules for $h$
heads. Choose one accepting path for each $(A,A^\top)$. Some schedule
occurs for at least the claimed fraction. Fix it for the rest of
the construction.

Use one vertex per epoch. Its incident bundles carry the entering and
leaving control state and head offsets, and the before and after
contents of the blocks in the three-block window around each starting
head. The number of blocks per epoch is constant. On a multitape
alphabet each block needs $O(N)$ bits. A table entry counts the legal
local paths of this epoch that realize these data and end in the
scheduled blocks. The first state is fixed and the last must accept.
Invalid boundary codes have entry zero.

Pass control and offsets to the next epoch. For each visited block,
connect successive uses in chronological order, including over epochs
that never visit it. Its contents are preserved in those intervening
epochs. Initialize work blocks with blank symbols and discard final
block contents by summation. For the input, attach a separate terminal
to each row of $A$ and each row of $B$. The latter are the column
indices of the desired transpose tensor. Each row intersects at most
two tape blocks even when the prefix shifts the alignment. A terminal
splits its exterior string into those pieces, and each initial block
joins its constant number of pieces with the fixed header and any
endmarkers. Unused pieces are summed against ones. These initial
tables cost $O(N+\lceil T/N\rceil)$ vertices, including blank blocks
visited outside the input. All connections preserve bounded degree
and width $O(N)$.

For fixed $A,B$ and a path following the schedule the boundary data
are unique. Conversely, for consistent boundary data any choices of
the local paths concatenate to a genuine path. Multiplication and
summation therefore give exactly the number of accepting paths with
that schedule. Soundness gives zero off the transpose relation. On
the selected inputs at least one path contributes, and its positive
coefficient survives over $\mathbb Q$. This proves the lemma.
\end{proof}

\begin{theorem}\label{thm:tape}
Assume the geometric estimate above. Every nondeterministic decider
for $\TrLang$ has worst-case time $\Omega(m_N\sqrt{\log m_N})$ at
every sufficiently large valid length $m_N$.
\end{theorem}
\begin{proof}
If $T>N^2\sqrt{\log N}$ the bound already holds. Otherwise the loss
in the exponent of the successful fraction in Lemma~\ref{lem:network}
is $O(N\sqrt{\log N})=o(N^2)$, within the allowance for any fixed
positive $\rho<1-1/\sqrt2$. Fixing a successful $A$ leaves precisely
$K(A,A^\top)e_{A^\top}$. The geometric bound now gives
\[
 N+\lceil T/N\rceil=\Omega(N\sqrt{\log N}),
 \qquad T=\Omega(N^2\sqrt{\log N}).
\]
Since $m_N=\Theta(N^2)$ the stated estimate follows.
\end{proof}

Combining the printed geometric theorem with Lemma~\ref{lem:upper}
gives, under Weak ETH,
\[
 \NT[n]\ne\coNT[n],\qquad
 \mathsf{BPTIME}[n]\not\subseteq\NT[o(n\sqrt{\log n})].
\]
The hard lengths here are $m_N$. The linear circuit in Lemma~\ref{lem:upper}
explains why the tape model matters.

\section{How the bound depends on the number of good arrays}\label{sec:audit}

Let $G$ be any nonempty set of $N\times N$ bit arrays. Suppose that a
rational network with $P$ vertices, degree at most $D$, and wire widths
$w_e\le DN$ has $2N$ distinct terminals $a_1,\ldots,a_N,b_1,\ldots,b_N$.
On each $A\in G$, pinning the row inputs must give
$c(A)e_{A^\top}$ with $c(A)\ne0$. There is no restriction on its output
outside $G$. Write $\beta=\log_2|G|/N^2$ and $z_+=\max(z,0)$.

\begin{theorem}\label{thm:density}
Under the characteristic-zero no-ports placement and the spectral wire
price comparison, there is a $c_D>0$ such that, for all sufficiently
large $N$ and every such $G$,
\[
       P\ \ge\ c_D N\bigl(1+(\beta-1/\sqrt2)_+\sqrt{\log N}\bigr).
\]
Both $c_D$ and the threshold for $N$ are independent of $G$.
\end{theorem}

\subsection{The three placement estimates}

Measure distances in the common normalization of the placement and set
$E=K_D P/\ln P$. Write $d_e$ for the distance between the ends of a
wire, $d_{ij}=d(a_i,b_j)$, and $o$ for the root of the placement.
The placement result gives the following three estimates simultaneously
\cite{partial,formalgeom}.
\begin{align}
 S=\sum_e w_e d_e &\le DN\sqrt{DP}\sqrt E, \label{eq:supply}\\
 \sum_{i,j}(d_{ij}-1)^2 &\le 12NE, \label{eq:terminals}\\
 \sum_v q_v d(v,o)&\le \frac{1}{\sqrt2}\sum_v q_v
               +\left(\sum_v q_v^2\right)^{1/2}\sqrt E
               \quad(q_v\ge0). \label{eq:root}
\end{align}
Here $P$ is the number of original vertices. The constant $K_D$ absorbs
the padding used to find a placement.

Two features of these estimates affect the exponent. To see the first,
pad the graph to $P'$ vertices with $P'=\Theta(P)$ as in the manuscript
and superimpose a bounded-degree scalar expander. The logarithmic
Poincar\'e obstruction under no ports~\cite{loggraph} supplies a
placement of the augmented graph with mean squared edge length
$O_D(1/\log P)$ and mean squared distance one between independent
uniform vertices $U,V$. Cauchy--Schwarz gives the supply bound.
Every real $1$-Lipschitz function on the placement has variance
$O_D(1/\log P)$ by the scalar expander inequality. This applies to
$d(U,t)$ for each fixed $t$ and to $g(V)=\mathbb E_U d(U,V)$.
By total variance, $\mu=\mathbb E d(U,V)=1+O_D(1/\log P)$.
Restricting $U$ to row terminals increases the mean squared deviation
from $g(t)$ by at most $P'/N$, uniformly in $t$. Restricting $V$ to
column terminals increases the mean squared deviation of $g(V)$ from
$\mu$ by the same factor. Adding the two deviations in $L^2$ incurs
the density loss $P'/N$ only once in the squared estimate.

The second feature is the coefficient of the root charge. Take $o$
to be the barycenter of the uniform vertex law in the complete
$\operatorname{CAT}(0)$ target. Its variance inequality gives
\[
 \mathbb E_U d(U,y)^2\ \ge\ \mathbb E_U d(U,o)^2+d(o,y)^2.
\]
Averaging $y=V$ yields $2\mathbb E d(U,o)^2\le1$. The mean distance
to $o$ is therefore at most $1/\sqrt2$. Apply the scalar variance
bound to $d(U,o)$ and Cauchy--Schwarz to the weighted deviations to
obtain~\eqref{eq:root} with a degree-dependent error constant.

\subsection{The entropy charged to the row and column terminals}

Draw $W$ uniformly from $G$, order its $N^2$ cells, and let
$h_{ij}=H(W_{ij}\mid W_{\mathrm{earlier}})$ in bits. Then
$0\le h_{ij}\le1$ and $\sum_{ij}h_{ij}=\log_2|G|$. The specialized
paired-terminal inequality is
\begin{equation}\label{eq:pins}
       \sum_{i,j}h_{ij}d_{ij}\ \le\ S+N\sum_i d(a_i,o).
\end{equation}
Indeed, a pin at $a_i$ specifies row $i$ and copies it to fresh external
indices there. Each bit is now displayed at both $a_i$ and $b_j$ with
the single joint coefficient $c(W)$. A row pin costs at most $N$ bits.
The other vertices have a single possible pin. These are exactly the
hypotheses of the paired-terminal entropy lemma~\cite{partial,formalent}.

We will also use a sharper version. Set $q_i=H(W_{i,*})$, the entropy
of the entire $i$th row. Then
\begin{equation}\label{eq:rowpins}
       \sum_{i,j}h_{ij}d_{ij}\ \le\ S+\sum_i q_i d(a_i,o).
\end{equation}
Here is a direct derivation for these pins. Take $r$ copies with $r$
divisible by $|G|$, and consider lists
in which every array of $G$ appears $r/|G|$ times. For cell $j$ in
the chosen ordering, and every prefix of the prescribed type, the
number of compatible next bit columns across the $r$ copies is
\[
 M_j=\prod_{u}\binom{r\Pr(W_{<j}=u)}
                        {r\Pr(W_{<j}=u,W_j=1)}.
\]
Each prefix extends to a full list by filling the positions of each
prefix word with its prescribed continuations. The binomial formula
gives $r^{-1}\log_2 M_j\to H(W_j\mid W_{<j})$. For instance,
if two cells range uniformly over $00,01,10$ and $r=3k$, the factors
are $\binom{3k}{k}$ and $\binom{2k}{k}$. Their combined rate is
$\log_2 3$, accounting for the correlation between the two cells.

For each $j$ color binary $r$-columns with $2^{\ell_j}$ colors, where
$\ell_j=\max(0,\lfloor\log_2(M_j/r^2)\rfloor)$. For $\ell_j>0$,
independent uniform colors miss a given color at a given prefix
with probability at most $\exp(-r^2)$. There are at most
$N^2 2^{r(N^2+1)}$ choices of cell, prefix and color. With $N$ fixed
and $r$ large, some coloring therefore offers every color at every
compatible prefix. A cell with one color already has this property.
Successively extending prefixes gives one full representative for
every tuple of colors.

At $o$, supply the sum of the representatives. The stream of row $i$
has fixed type, so its number of possible values is at most $2^{rq_i}$.
Indeed, each such stream has probability $2^{-rq_i}$ in $r$ independent
draws from the row marginal. A $\lceil rq_i\rceil$-bit wire to $a_i$
therefore suffices, with the index decoded by its local map.
Their common scalar $\prod_{W\in G}c(W)^{r/|G|}$ is nonzero.
Divide the root tensor by this scalar.
The full set of row streams distinguishes the representatives. Apply
the corresponding cell's color map at both ends. The resulting tensor
is a product of independent equality wires, one $\ell_j$-bit wire for
each cell, each joint label appearing with coefficient one. Spectral
prices add under products and cannot increase under these local maps.
The $r$ copies cost $rS$ and supplying the rows costs at most
$\sum_i\lceil rq_i\rceil d(a_i,o)$. Dividing by $r$ and taking the
limit proves~\eqref{eq:rowpins}, hence also~\eqref{eq:pins}.
The auxiliary root and its wires enter only this
price comparison. The degree bound concerns the original network.

\begin{proof}[Proof of Theorem~\ref{thm:density}]
Put $H=\log_2|G|$ and $R=\sum_i q_i$. By~\eqref{eq:terminals} and
Cauchy--Schwarz, replacing all $d_{ij}$ by $1$ on the left of
\eqref{eq:rowpins} incurs error at most $N\sqrt{12NE}$.
In~\eqref{eq:root} use weight $q_i\le N$ at row terminal $a_i$ and
zero elsewhere. The root charge is at most $R/\sqrt2+N\sqrt{NE}$.
By~\eqref{eq:supply},
\[
  \frac{H-R/\sqrt2}{N^2}
     \le C_D\left(\frac{P}{N\sqrt{\log P}}
                   +\sqrt{\frac{P}{N\log P}}\right).
\]
Since $P\ge2N$, the second term is at most the first and
$\log P\ge\log N$. We have proved the stronger estimate
\begin{equation}\label{eq:rowgain}
 H\ \le\ \frac{R}{\sqrt2}+O_D\!\left(\frac{NP}{\sqrt{\log N}}\right).
\end{equation}
Using $R\le N^2$ and $P\ge2N$ proves the theorem.
\end{proof}

For $|G|\ge 2^{(1-\rho)N^2}$ with fixed
$\rho<1-1/\sqrt2$, Theorem~\ref{thm:density} recovers the printed
partial-transposition theorem. At $\beta=1/\sqrt2$ this argument
gives only the linear bound from the terminals.

\begin{corollary}[Fraction of true pairs with a short accepting path]
\label{cor:coverage}
Assume the hypotheses of Theorem~\ref{thm:density}. Let a fixed
nondeterministic machine be sound on transpose pairs, meaning that
it never accepts a false pair. For $T_N=o(N^2\sqrt{\log N})$, let
$\pi_N$ be the fraction of matrices $A$ for which some path accepts
$(A,A^\top)$ within $T_N$ steps. Then
\[
                 \pi_N\ \le\ 2^{-(1-1/\sqrt2-o(1))N^2}.
\]
The machine may run longer on other paths or fail to accept some
true pairs altogether.
\end{corollary}
\begin{proof}
Assume $\pi_N>0$ and truncate at $T_N$. The schedule argument of
Lemma~\ref{lem:network} retains a set of size at least
$2^{N^2}\pi_N\,3^{-h\lceil T_N/N\rceil}$ for $h$ tape heads.
Its network has $P=O_M(N+\lceil T_N/N\rceil)=o(N\sqrt{\log N})$.
Theorem~\ref{thm:density} gives $\log_2|G|/N^2\le1/\sqrt2+o(1)$.
The schedule costs $O_M(1+T_N/N)=o(N^2)$ bits in the logarithm.
Rearranging gives the claim. The case $\pi_N=0$ is immediate.
\end{proof}

\begin{corollary}[Prescribed sets of rows]\label{cor:product}
For each $N$, let $\mathcal R=\prod_{i=1}^N\mathcal R_i$, where
$\varnothing\ne\mathcal R_i\subseteq\{0,1\}^N$, and put
$H_0=\log_2|\mathcal R|$. Under the same hypotheses, the fraction
$\pi_{\mathcal R}$ of $A\in\mathcal R$ with an accepting path by time
$T$ satisfies, for all sufficiently large $N$,
\[
 \pi_{\mathcal R}\ \le\
 2^{-(1-1/\sqrt2)H_0+
        O_M\bigl((N^2+T)/\sqrt{\log N}+1+T/N\bigr)}.
\]
In particular, if $H_0\ge\delta N^2$ for a fixed $\delta>0$ and
$T=o(N^2\sqrt{\log N})$, the relative fraction is at most
$2^{-(1-1/\sqrt2-o(1))H_0}$.
\end{corollary}
\begin{proof}
When $\pi_{\mathcal R}>0$, choose one timely accepting path for each
accepted array in $\mathcal R$. A common schedule retains
$H=\log_2|G|\ge H_0+\log_2\pi_{\mathcal R}-O_M(1+T/N)$ bits.
Conditioning on that schedule can correlate the rows. Nevertheless,
the $i$th row still belongs to $\mathcal R_i$, so
$R\le\sum_i\log_2|\mathcal R_i|=H_0$. Apply~\eqref{eq:rowgain} with
$P=O_M(N+\lceil T/N\rceil)$ and rearrange. The same proof only needs
soundness for pairs $(A,B)$ with $A\in\mathcal R$ and $B$ arbitrary,
since every pinned $A$ lies in $\mathcal R$.
\end{proof}
For example, each row may range over a different codebook of
$2^{\lfloor N/10\rfloor}$ words. Even accepting every true pair in
this family covers only a $2^{-0.9N^2+O(N)}$ fraction of all matrices.
The absolute density theorem alone gives no superlinear bound at that
density. Corollary~\ref{cor:product} still forces
$\Omega_M(N^2\sqrt{\log N})$ time for complete coverage of these
codebooks.

\subsection{What the released proofs cover}

The printed placement proof uses Weak ETH to exclude ports.
This implication involves the variable count: enumerating the contraction
of a sparse formula does not by itself give a power saving for a dense
one. In the source's sunflower branching, $r=\lceil v^\alpha\rceil$,
the tree has at most $2^{O((v/r)\log v)}$ leaves, and each leaf has
a bounded-degree network of volume $O(vr^4)$~\cite{expansion}.
For a contraction exponent $1-\eta$, choose fixed $\alpha>0$ with
$(1+4\alpha)(1-\eta)<1$. Both terms in the running-time exponent now
have a power saving in $v$.

The logarithmic factor in the placement requires a quantitative
obstruction. In the proof of~\cite{loggraph}, suppose a $d$-regular
$s$-vertex graph has a Poincar\'e bound $\lambda$ throughout the class
of spectral spaces. Starting from a regular graph of degree $2d$ with
bound $\gamma$, average the first $t$ powers of its lazy walk. The
Hadamard walk-smoothing lemma improves the bound for this average to
$O(1+\gamma/t)$. For even $t$ with $(2d+1)^t\le s/2$, each row of
the averaged kernel has at most $s/2$ entries in its support. Round
each entry $p$ to $\lceil sp/2\rceil$ arcs, then complete each row
to $s$ arcs by adding loops. Symmetry is preserved and the resulting
transition probabilities dominate half the original ones. Replace
each old vertex by a copy of the $s$-vertex graph and match vertices
of these copies according to the rounded arc counts. Comparing within
copies and between their barycenters gives the recurrence
$\gamma'\le B\max(1,\lambda)(1+\gamma/t)$ for an absolute $B$.
If $t\ge2B\max(1,\lambda)$, iteration gives arbitrarily large graphs
of the same bounded degree with uniformly bounded Poincar\'e constants.
That is the ports property. Choosing the largest admissible even $t$
explains the $\Omega_d(\log s)$ obstruction under no ports. Realizability
is used later, when the resulting metric distances serve as tensor prices.

There is also a route through the precise declarations in the pinned
Lean sources. In characteristic zero the formal alternative uses
$\overline{\mathbb Q}$. Embedding our rational coefficients in this field
preserves every contraction, zero coefficient, and nonzero scalar.
The theorem \texttt{graph\_geometry} gives
\eqref{eq:supply}--\eqref{eq:root} under \texttt{PriceZero}, using one
target, one positive scale and the original vertex count~\cite{formalgeom}.
Finite realizability applies jointly to the vertices and root. The common
scale multiplies both sides of the price comparison.
In \texttt{paired\_entropy\_external\_root}, take the finite law to be
uniform on $G$, the separating labels to be its individual bits, the two
terminal maps to be their row and column, and the local maps to be
the row-pin construction above~\cite{formalent}. The chain-rule
entropies in that declaration are divided by $\ln2$, hence give bits.
Finally, \texttt{price\_of\_variable\_weakETH} gives \texttt{PriceZero}
under the variable-count hypothesis, and \texttt{middle\_count\_bridge}
gives the decision algorithm under its opposite, \texttt{PortZero}~\cite{formaleth}.
This states both branches using the same formal definition of ports.
The released comparison endpoint concerns multiplication on tapes~\cite{scope}.
The specialization to this decision language and the density calculation
above are paper proofs. This audit compares the declarations and their
mathematical hypotheses. A fresh Lean build and Comparator run are
outside the checks reported here.

\section{A circuit alternative without Weak ETH}\label{sec:cases}
In the terminology of the tensor manuscripts, ports are an expansion
property of the spectral spaces of a characteristic. In addition to
the placement and price hypotheses of Theorem~\ref{thm:density}, assume the enumerable
contraction estimate in characteristic zero with ports~\cite{expansion}.
For sparse 3-SAT it gives a deterministic $V^{O(1)}2^{O(V^{1-\eta})}$
algorithm for some fixed $0<\eta<1$. Indeed a width-three formula on
$V$ declared variables and at most $5V$ clauses has a closed counting
network of volume $O(V)$: use a clause predicate per clause and copy
each repeated variable by a tree of ternary equality tables. Its integer
contraction counts assignments (unused variables may be omitted for
decision). Entries are bits and degree is bounded, so the enumerable
contraction estimate applies directly for large volumes.

\Needspace{6\baselineskip}
\begin{corollary}
Under these two geometric hypotheses at least one of the following holds:
\begin{enumerate}
\item $\TrLang$ obeys the lower bound of Theorem~\ref{thm:tape}.
\item For some $\varepsilon,c>0$ there is a language $H\in\ENP$ with
      $C_{B_2}(H_n)\ge cn^{1+\varepsilon}$ at every sufficiently large $n$.
\end{enumerate}
\end{corollary}
\begin{proof}
Without ports, apply Theorem~\ref{thm:tape}. With ports, use the sparse
algorithm above and the inverse-rate theorem of Result~02~\cite{inverse}
with $\alpha=1-\eta$. A gate budget of order $n^{1/\alpha}$ fits its
verifier clock, giving $\varepsilon=\eta/(1-\eta)$.
\end{proof}
Thus a verifier for this particular transpose language in
$o(n\sqrt{\log n})$ time would select the circuit alternative.

The choice of the tape problem matters here. A disjunction using only
$\NT[n]\ne\coNT[n]$ has an elementary proof with a stronger circuit
alternative. Equality of these linear-time classes would give
$\mathsf{NP}=\mathsf{coNP}$ by polynomial padding. Put $M=2^n$.
For sufficiently small fixed $c>0$, counting circuit descriptions gives
an $M$-bit truth table requiring at least $cM/n$ gates. The question
whether a given prefix extends to such a table is in $\Sigma_2^{\mathsf P}$
with $M$ in unary: guess a completion and check that every smaller
circuit disagrees somewhere by scanning all $M$ entries. Under the
collapse this predicate lies in $\mathsf{NP}$. Choosing the first
extendible bit at each step takes polynomial time in $M$ with an NP
oracle, giving the same hard table for all $n$-bit addresses.
For randomized linear time, padding a containment in nondeterministic
polynomial time gives $\mathsf{BPP}\subseteq\mathsf{NP}$ for total languages.
The near-maximum hardness theorem of Aaronson et al. uses a simulation
of a particular promise-AM approximate-counting problem~\cite{promise}.

The remaining imported steps are the no-ports placement in a realizable
spectral target and the ports contraction algorithm. The derivation in
Section~\ref{sec:audit} checks how the former gives the verifier bound
and where the SAT hypothesis enters. Historical priority for the
verifier consequences remains open.

\paragraph{Acknowledgments.}
AI was used to produce the results, proofs, and text of this note.

\begin{thebibliography}{9}
\small
\bibitem{partial} OpenAI. \emph{The logarithmic exponent of integer
multiplication on tapes}, 24 November 2026. Theorem on partial
transposition and lemma on short edges and distant terminals.
\href{https://github.com/openai/math/blob/3344da59991ed94f9ba711217b29c1d760badfa6/preprints/The-logarithmic-exponent-of-integer-multiplication-on-tapes-November-24-2026/build/lower.tex}{Pinned proof}.
\bibitem{expansion} OpenAI. Same manuscript, theorem on expansion and
3-SAT, including its statement of enumerable contraction.
\href{https://github.com/openai/math/blob/3344da59991ed94f9ba711217b29c1d760badfa6/preprints/The-logarithmic-exponent-of-integer-multiplication-on-tapes-November-24-2026/build/graphs.tex}{Pinned source}.
\bibitem{loggraph} OpenAI. \emph{Circuit lower bounds for Hamiltonian
cycle and the permanent}, 24 November 2026.
\href{https://github.com/openai/math/blob/3344da59991ed94f9ba711217b29c1d760badfa6/preprints/Circuit-lower-bounds-for-Hamiltonian-cycle-and-the-permanent-November-24-2026/build/entropy.tex}{Logarithmic Poincar\'e obstruction}
and \href{https://github.com/openai/math/blob/3344da59991ed94f9ba711217b29c1d760badfa6/preprints/Circuit-lower-bounds-for-Hamiltonian-cycle-and-the-permanent-November-24-2026/build/compress.tex}{enumerable contraction under ports}.
\bibitem{formalgeom} OpenAI. \emph{MultitapeWeakETH}, Lean sources at
revision \texttt{3344da5}. \href{https://github.com/openai/math/blob/3344da59991ed94f9ba711217b29c1d760badfa6/lean/OAI/Computability/MultitapeWeakETH/Padding.lean}{Graph geometry on the original vertex set}.
\bibitem{formalent} OpenAI. Same revision.
\href{https://github.com/openai/math/blob/3344da59991ed94f9ba711217b29c1d760badfa6/lean/OAI/Computability/MultitapeWeakETH/ExternalRoot.lean}{Entropy with an external root} and
\href{https://github.com/openai/math/blob/3344da59991ed94f9ba711217b29c1d760badfa6/lean/OAI/Computability/MultitapeWeakETH/BitRates.lean}{sum and range of conditional bit entropies}.
\bibitem{formaleth} OpenAI. Same revision.
\href{https://github.com/openai/math/blob/3344da59991ed94f9ba711217b29c1d760badfa6/lean/OAI/Computability/MultitapeWeakETH/Price.lean}{Characteristic-zero alternatives} and
\href{https://github.com/openai/math/blob/3344da59991ed94f9ba711217b29c1d760badfa6/lean/OAI/Computability/MultitapeWeakETH/MiddleDecision.lean}{variable-count decision theorem}.
\bibitem{scope} OpenAI. Same revision.
\href{https://github.com/openai/math/blob/3344da59991ed94f9ba711217b29c1d760badfa6/lean/docs/379.md}{Formalization scope for integer multiplication on tapes}.
\bibitem{inverse} C.~Cruz. \emph{Almost-everywhere circuit lower bounds
for $\ENP$ from sparse 3-SAT algorithms}. Result~02, section on inverse
rates and unbounded gates, revised November 2026.
\bibitem{promise} S.~Aaronson, B.~Ayd\i nl\i o\u{g}lu, H.~Buhrman,
J.~Hitchcock and D.~van Melkebeek. \emph{A note on exponential
circuit lower bounds from derandomizing Arthur--Merlin games},
Theorem~3. \href{https://www.cs.uwyo.edu/~jhitchco/papers/neclb.pdf}{Paper}.
\end{thebibliography}
\end{document}
